\documentclass[11pt]{article}
\usepackage[margin=1in]{geometry}
\usepackage{amsmath,amssymb,amsthm,mathtools}
\usepackage{hyperref}
\usepackage{enumitem}
\usepackage{booktabs}
\usepackage{xcolor}

\theoremstyle{plain}
\newtheorem{proposition}{Proposition}[section]
\theoremstyle{definition}
\newtheorem{definition}{Definition}[section]
\newtheorem{conjecture}{Conjecture}[section]
\theoremstyle{remark}
\newtheorem{remark}{Remark}[section]

\hypersetup{colorlinks=true, linkcolor=blue, citecolor=blue, urlcolor=blue, pdfusetitle}

\title{Entanglement-Entropy Sourcing of Gravity: Dimensional Repair,\\Covariant Completion, and Experimental Outlook}
\author{Kevin Monette}
\date{March 19, 2026}

\begin{document}
\maketitle

\begin{abstract}
We develop a dimensionally consistent phenomenological extension of thermodynamic gravity
in which a \emph{protocol-certified} internal entanglement entropy density of a quantum
device acts as an additional gravitational source.  The central repair over earlier drafts
is twofold: a reference energy scale $E_*$ (J/bit) converts dimensionless entanglement
counts to an energy density, and the source field is redefined from an ambiguous
``internal entropy'' to $S_{\mathrm{cert}}(x)$ --- entanglement generated by a specified
control sequence, certified by a witness, and varied between experimental branches.
This operational definition, grounded in the theory of accessible entanglement under
superselection rules, excludes uncontrolled background correlations.  The Cooper-pair
constraint is resolved not by asserting that BCS entanglement is absent --- it is not ---
but by a \textbf{Cancellation Proposition}: because the substrate state is identical
in both experimental branches, background entanglement cancels exactly in the differential
observable $\Delta S_{\mathrm{cert}}$.  An explicit open conjecture (OP-new) is stated
for the residual contributions.  Under the Planck-scale benchmark $E_* = E_P$, all
observables reduce to $\Lambda_{\mathrm{eff}} = \tilde{\kappa}E_*\alpha_{\mathrm{screen}}$,
with a predicted unsuppressed acceleration of $\Delta a = 3.6\times10^{-9}$\,m/s$^2$
for 100 protocol-certified entangled bits at $100\,\mu$m, within reach of cryogenic
SQUID-based force sensors.
\end{abstract}

\section{Purpose and Scope}

This document is a dimensional-consistency revision of the entanglement--gravity program. Its goal is not to select a unique microscopic origin for the coupling, but to isolate the minimal scaling assumptions required for a well-posed phenomenological theory.

Earlier drafts used a coupling prefactor proportional to $c^4/(8\pi G k_B \ln 2)$ multiplying an entanglement entropy density measured in bits per unit volume. In SI units, that expression does not yield an energy density, so the resulting stress--energy tensor and derived accelerations were not dimensionally valid. The present revision fixes that problem by introducing a reference energy scale $E_*$ and rebuilding the source term around it.

\section{Operational Definition and Units}

\subsection{The dimensional problem}

The Einstein equations
\begin{equation}
G_{\mu\nu} = \frac{8\pi G}{c^4} T_{\mu\nu}
\end{equation}
require $[T_{\mu\nu}] = \mathrm{J\,m^{-3}}$.  Writing $[S_{\mathrm{ent}}] = \mathrm{bit\,m^{-3}}$
(bit dimensionless) and using the original coupling $\alpha = \tilde{\kappa}c^4/(8\pi Gk_B\ln 2)$
gives $[\alpha S_{\mathrm{ent}}] = \mathrm{K\,m^{-4}}$, not an energy density.  A reference energy
scale $E_*$ (J/bit) is the minimal repair.

\subsection{Protocol-certified entanglement: \texorpdfstring{$S_{\mathrm{cert}}$}{Scert}}
\label{sec:Scert}

The earlier notation $S_{\mathrm{cert}}$ was ambiguous: it could be interpreted as the von Neumann
entropy of any reduced density matrix, including thermal and BCS correlations in the device
substrate.  We replace it with a precisely defined operational quantity.

\begin{definition}[Protocol-certified entanglement density]
Let $D$ be a quantum device, $P$ a specified gate sequence
acting on the device degrees of freedom, $D = A\cup B$ a bipartition pre-registered before
$P$ is applied, and $\mathcal{W}$ an entanglement witness verified after $P$.
The \textbf{protocol-certified entanglement density} is
\begin{equation}
S_{\mathrm{cert}}(x)
= \frac{-k_B\,\mathrm{Tr}_A\!\left[\rho_A^P(x)\ln\rho_A^P(x)\right]}{k_B\ln 2\,\Delta V},
\quad
\rho_A^P = \mathrm{Tr}_B\!\left[P\rho_D P^\dagger\right],
\label{eq:Scert}
\end{equation}
where $P\rho_D P^\dagger$ is the device state after protocol $P$ and $\Delta V$ is the cell
volume.  We require that $\mathcal{W}$ certifies $S_{\mathrm{cert}} > 0$ to within
experimental error before a run is accepted.
\end{definition}

This replaces the generic bipartite von Neumann entropy used in earlier drafts.  The definition
has three operationally necessary properties:
(i)~$S_{\mathrm{cert}} \geq 0$;
(ii)~$S_{\mathrm{cert}} = 0$ if no entangling protocol was applied;
(iii)~$S_{\mathrm{cert}}$ is certified by a positive entanglement witness, so it cannot be
mimicked by classical correlations.


\section{Minimal Dimensional Fix}

\subsection{Reference energy scale}

We introduce a reference energy scale $E_*$ with units of energy per bit and define the entanglement energy density by
\begin{equation}
\rho_{\mathrm{ent}}(x) = E_*\,S_{\mathrm{ent}}(x).
\end{equation}
Then
\begin{equation}
[\rho_{\mathrm{ent}}] = [E_*][S_{\mathrm{ent}}] = \mathrm{J\,m^{-3}},
\end{equation}
as required.

An equivalent parameterization is
\begin{equation}
\rho_{\mathrm{ent}}(x) = \lambda_* k_B T_* s_{\mathrm{ent}}(x),
\end{equation}
where $s_{\mathrm{ent}}$ is entropy density in thermodynamic units, $T_*$ is an effective temperature scale, and $\lambda_*$ is dimensionless. In this white paper we use the simpler single-scale notation $E_*$.

\subsection{Revised source term}

The entanglement contribution to the stress--energy tensor is rewritten as
\begin{equation}
T^{\mathrm{ent}}_{\mu\nu}(x) = \tilde{\kappa}\,\rho_{\mathrm{ent}}(x) g_{\mu\nu} + \text{(gradient terms)},
\end{equation}
with
\begin{equation}
\rho_{\mathrm{ent}}(x) = E_* S_{\mathrm{ent}}(x).
\end{equation}
Thus all previous appearances of $\alpha S_{\mathrm{ent}}$ are replaced by $\tilde{\kappa} E_* S_{\mathrm{ent}}$.

If desired, one may later relate $E_*$ to a microscopic scale such as
\begin{equation}
E_* = \eta k_B T_*,
\end{equation}
with $\eta$ dimensionless and $T_*$ an effective Unruh or horizon temperature. The present framework leaves the origin of $E_*$ open.

\section{Revised Effective Action and Fifth-Force Equation}

\subsection{External-source action}

Let
\begin{equation}
\phi(x) \equiv S_{\mathrm{ent}}(x).
\end{equation}
We retain the external-source interpretation: $\phi$ is prescribed by the device quantum state and is not varied as an independent field. The entanglement-sector Lagrangian density is written as
\begin{equation}
\mathcal{L}_{\mathrm{ent}} = -\frac{\beta}{2} g^{\mu\nu}(\nabla_\mu\phi)(\nabla_\nu\phi) - \lambda\phi,
\end{equation}
with
\begin{equation}
\lambda = \tilde{\kappa} E_*.
\end{equation}
The full action is therefore
\begin{equation}
S[g;\phi] = \int d^4x\,\sqrt{-g}\left[\frac{c^4}{16\pi G}R - \frac{\beta}{2} g^{\mu\nu}(\nabla_\mu\phi)(\nabla_\nu\phi) - \tilde{\kappa} E_*\phi\right].
\end{equation}

Varying with respect to $g_{\mu\nu}$ yields
\begin{equation}
T^{\mathrm{ent}}_{\mu\nu} = \beta\left(\nabla_\mu\phi\,\nabla_\nu\phi - \frac{1}{2}g_{\mu\nu}(\nabla\phi)^2\right) + \tilde{\kappa} E_*\phi\, g_{\mu\nu}.
\end{equation}
The dimensions are correct provided $\beta$ carries the dimensions needed to make the gradient term an energy density. In slowly varying configurations, the gradient term is subleading.

\subsection{Matter nonconservation}

Total conservation requires
\begin{equation}
\nabla^\mu T^{\mathrm{total}}_{\mu\nu} = 0,
\end{equation}
so that
\begin{equation}
\nabla^\mu T^{\mathrm{matter}}_{\mu\nu} = -\nabla^\mu T^{\mathrm{ent}}_{\mu\nu}.
\end{equation}
Substituting the revised entanglement tensor gives
\begin{equation}
\nabla^\mu T^{\mathrm{matter}}_{\mu\nu} = -\left(\tilde{\kappa} E_* + \beta\Box\phi\right)\nabla_\nu\phi + \beta\nabla_\nu\left(\frac{1}{2}(\nabla\phi)^2\right).
\end{equation}
In the quasi-static regime with modest gradients and $|\beta\Box\phi| \ll |\tilde{\kappa}E_*|$, this reduces to
\begin{equation}
\nabla^\mu T^{\mathrm{matter}}_{\mu\nu} \approx -\tilde{\kappa}E_*\,\nabla_\nu\phi(x) = -\tilde{\kappa}E_*\,\nabla_\nu S_{\mathrm{ent}}(x).
\end{equation}
This is the corrected fifth-force law. Its units are now those of force density, as required.

\section{Newtonian Limit and Revised Observable}

\subsection{Modified Poisson equation}

In the Newtonian limit,
\begin{equation}
\nabla^2 \Phi = 4\pi G\left(\rho_{\mathrm{matter}} + \rho_{\mathrm{ent}}\right),
\end{equation}
with
\begin{equation}
\rho_{\mathrm{ent}}(\mathbf{x}) = \tilde{\kappa} E_* S_{\mathrm{ent}}(\mathbf{x}).
\end{equation}
Under spherical symmetry,
\begin{equation}
\Delta a(R) = -\frac{d\Phi_{\mathrm{ent}}}{dR} = \frac{G}{R^2}M_{\mathrm{ent}}(<R),
\end{equation}
where
\begin{equation}
M_{\mathrm{ent}}(<R) = 4\pi\int_0^R \rho_{\mathrm{ent}}(r)r^2\,dr.
\end{equation}
For bounded sources, $M_{\mathrm{ent}}(<R)$ saturates outside the source and the correction scales as $1/R^2$.

\subsection{Observable combination}

Experiments constrain the product
\begin{equation}
\Lambda_{\mathrm{eff}} \equiv \tilde{\kappa}\,E_*\,\alpha_{\mathrm{screen}},
\quad [\Lambda_{\mathrm{eff}}] = \mathrm{J/bit},
\end{equation}
not $\tilde{\kappa}$ or $E_*$ individually. In branch-comparison experiments the
useful observational combination is
\begin{equation}
\mathcal{B}
\equiv |\Lambda_{\mathrm{eff}}|\,\frac{\Delta S_{\mathrm{ent}}^{\mathrm{tot}}}{R_{\mathrm{eff}}^2},
\quad [\mathcal{B}] = \mathrm{J\,m^{-3}},
\end{equation}
where $\Delta S_{\mathrm{ent}}^{\mathrm{tot}}$ is the differential internal entanglement content between
experimental branches and $R_{\mathrm{eff}}$ is the characteristic geometric scale. A null result
at sensitivity $\Delta a_{\max}$ constrains
\begin{equation}
|\Lambda_{\mathrm{eff}}| < \frac{\Delta a_{\max}\,R_{\mathrm{eff}}^2}{G\,\Delta S_{\mathrm{ent}}^{\mathrm{tot}}}.
\end{equation}
Mapping this bound to $|\tilde{\kappa}|$ alone requires an independent assumption about
$E_*$; without such an assumption, only $|\Lambda_{\mathrm{eff}}|$ is constrained.

\section{Implications for OP8 and OP10}

\subsection{Fifth-force bounds}

The corrected acceleration outside a bounded source of radius $R_0$ is
\begin{equation}
\Delta a(R) = \frac{G\,|\Lambda_{\mathrm{eff}}|\,\Delta S_{\mathrm{ent}}^{\mathrm{tot}}}{R^2},
\quad R > R_0.
\end{equation}
A null force bound $\Delta a_{\max}$ constrains the observable product
\begin{equation}
|\Lambda_{\mathrm{eff}}|
= |\tilde{\kappa}|\,E_*\,\alpha_{\mathrm{screen}}
\lesssim \frac{\Delta a_{\max}\,R_{\mathrm{eff}}^2}{G\,\Delta S_{\mathrm{ent}}^{\mathrm{tot}}}.
\label{eq:Lambda-bound}
\end{equation}

\paragraph{Benchmark (unsuppressed signal).}
For the Planck-scale hypothesis $E_* = E_P = 1.96\times10^9$\,J and $|\tilde{\kappa}| = 1/4$, the
effective mass per bit is $m_{\mathrm{bit}} = |\tilde{\kappa}|E_P/c^2 = 5.44\,\mu\mathrm{g/bit}$.
For $\Delta S_{\mathrm{ent}} = 100$\,bits at $R = 100\,\mu\mathrm{m}$ with $\alpha_{\mathrm{screen}} = 1$:
\begin{equation}
\Delta a
= \frac{G \times 100 \times 5.44\times10^{-9}\,\mathrm{kg}}{(10^{-4}\,\mathrm{m})^2}
= 3.6\times10^{-9}\,\mathrm{m/s^2}.
\end{equation}
Cryogenic SQUID-based force sensors achieve $\sim 10^{-10}$\,m/s$^2$/\hspace{0.5pt}$\sqrt{\mathrm{Hz}}$,
placing this signal above the noise floor by a factor of $\sim\!36$ at 1\,Hz bandwidth.

\paragraph{Critical constraint.}
If $E_* \approx E_P$, macroscopic superconductors containing $\sim\!10^{20}$ Cooper pairs would
carry an effective entanglement mass of order $10^{11}$\,kg per gram of material, producing
enormous gravitational anomalies not observed. The framework requires that the operational
definition of $S_{\mathrm{cert}}$ excludes uncontrolled thermal correlations (including BCS
Cooper-pair entanglement) for which no controlled bipartition protocol was applied. Verifying
this claim rigorously is an open problem; see Section~\ref{sec:open-problems}.

\subsection{RG interpretation}

With $E_*$ explicit, the natural renormalization-group ansatz becomes
\begin{equation}
\tilde{\kappa}_{\mathrm{lab}}E_{*,\mathrm{lab}} = \tilde{\kappa}_{\mathrm{UV}}E_{*,\mathrm{UV}}\left(\frac{E_{\mathrm{lab}}}{E_P}\right)^\gamma,
\end{equation}
where $E_P$ is the Planck energy and $\gamma$ is an effective scaling exponent. Experiments therefore constrain the infrared product $\tilde{\kappa}E_*$ rather than $\tilde{\kappa}$ alone.


\section{Operational Sourcing, Cancellation, and Open Problems}
\label{sec:open-problems}

\subsection{Why $S_{\mathrm{cert}}$ is the right object}
\label{sec:Scert-justification}

The definition of $S_{\mathrm{cert}}$ in Section~\ref{sec:Scert} may appear
restrictive.  This subsection explains why it is the \emph{only} definition
consistent with the goals of the framework.

The gravitational source must satisfy three physical requirements:
\begin{enumerate}[label=(\roman*),noitemsep]
  \item \textbf{Branch-modulation:} it must differ between the high-entanglement branch
    A and the reference branch B, otherwise no differential signal exists.
  \item \textbf{Controllability:} it must be adjustable by the experimenter, not fixed
    by the material.
  \item \textbf{Certifiability:} it must be verifiable that entanglement, not
    classical correlations, is responsible for any signal.
\end{enumerate}
The von Neumann entropy of an arbitrary reduced density matrix fails (i): BCS ground
states, thermal states of any system in equilibrium, and nuclear spins all contribute
a background $S_{\mathrm{vN}}$ that is fixed by the material and identical in both branches.
$S_{\mathrm{cert}}$ satisfies all three requirements by construction.

This is not merely a convenient definition --- it is the operationally accessible
entanglement in the sense of Wiseman and Vaccaro~\cite{Wiseman2003}: entanglement
that can be extracted and used as a resource under the device's allowed operations
(local unitary gates on the qubit register, projective readout, classical communication).
Background correlations that are not modifiable by the protocol are not resources
under this definition, and are therefore not sources in this framework.

\subsection{Cancellation theorem for background correlations}
\label{sec:cancel}

\begin{proposition}[Background cancellation]
\label{prop:cancel}
Let the total system be $T = D \cup B_{\mathrm{g}}$, where $D$ is the quantum device
and $B_{\mathrm{g}}$ is the background (substrate, Cooper pairs, thermal phonons).
Assume:
\begin{enumerate}[label=(\roman*),noitemsep]
  \item \textbf{Product initial state:}
    $\rho_T^{(0)} = \rho_D^{(0)} \otimes \rho_{B_{\mathrm{g}}}^{(0)}$\,.
  \item \textbf{Protocol locality:} the branch protocol acts on $D$ only:
    $P^{A,B} = U_D^{A,B} \otimes \mathbf{I}_{B_{\mathrm{g}}}$.
  \item \textbf{Bipartition locality:} the bipartition $\mathcal{A}$ is pre-registered
    on $D$ only, so the witness $\mathcal{W}$ measures correlations within $D$.
\end{enumerate}
Then the differential protocol-certified entanglement of the background vanishes exactly:
\begin{equation}
  \Delta S_{\mathrm{cert}}^{B_{\mathrm{g}}}
  \equiv S_{\mathrm{cert}}\!\left(\rho_{B_{\mathrm{g}}}; P^A, \mathcal{A}\right)
       - S_{\mathrm{cert}}\!\left(\rho_{B_{\mathrm{g}}}; P^B, \mathcal{A}\right)
  = 0.
  \label{eq:cancel}
\end{equation}
The measured signal depends only on $\Delta S_{\mathrm{cert}}^D = S_{\mathrm{cert}}^A
- S_{\mathrm{cert}}^B$ for the device qubits.
\end{proposition}

\begin{proof}
Under assumption~(ii), $P^{A,B}$ acts as the identity on $B_{\mathrm{g}}$.
Therefore $P^{A,B}\rho_{B_{\mathrm{g}}}P^{A,B\dagger} = \rho_{B_{\mathrm{g}}}$
for both branches.  Under assumption~(iii), the witness $\mathcal{W}$ measures
correlations across the bipartition $\mathcal{A}$ of $D$; it has no action on
$B_{\mathrm{g}}$.  Consequently
$S_{\mathrm{cert}}(\rho_{B_{\mathrm{g}}}; P^A) = S_{\mathrm{cert}}(\rho_{B_{\mathrm{g}}}; P^B) = 0$,
and Eq.~\eqref{eq:cancel} follows.
\end{proof}

\paragraph{Physical meaning.}
The Cooper-pair constraint is resolved not by claiming BCS entanglement is absent ---
it is not --- but by the observation that BCS correlations are \emph{branch-independent}.
The substrate state $\rho_{B_{\mathrm{g}}}$ is identical in branches A and B: both
branches run the same physical setup (same temperature, same chip, same equilibrium state)
with the only difference being the CZ gate phase applied to the qubit register.  A state
that is identical in both branches contributes zero to the differential observable
$\Delta S_{\mathrm{cert}}$ and zero to the gravitational signal.

\paragraph{Validity of assumption~(i).}
For superconducting transmon qubits, the device-substrate coupling is phononic
($J_{\mathrm{ph}} \lesssim 1$\,kHz) relative to qubit transition frequencies
($\omega/2\pi \sim 5$\,GHz), giving a ratio $J_{\mathrm{ph}}/\omega \sim 10^{-7}$.
The initial state is therefore a product state to excellent approximation.
Gate-induced phonon backaction creates device-substrate entanglement as a \emph{systematic},
which must be matched between branches (identical gate sequences on branch B); it does
not constitute $S_{\mathrm{cert}}$ under Definition~\ref{def:Scert}.

\subsection{OP-new: thermal equilibrium and $S_{\mathrm{cert}}$}
\label{sec:OP-new}

Proposition~\ref{prop:cancel} resolves the Cooper-pair constraint for differential
observables.  A residual question remains about the \emph{absolute} gravitational
source from any system in thermal equilibrium, independent of branch comparison.

\begin{conjecture}[OP-new: Thermal state conjecture]
\label{conj:thermal}
For any quantum system in a Gibbs state $\rho_\beta \propto e^{-\beta H}$, subject to
the operational constraints (local controls, charge/particle-number superselection,
fixed bipartition) of a realistic laboratory device, the protocol-certified entanglement
satisfies
\begin{equation}
  S_{\mathrm{cert}}(\rho_\beta; P_{\mathrm{id}}, \mathcal{A}, \mathcal{W})
  = 0,
  \label{eq:thermal-conj}
\end{equation}
where $P_{\mathrm{id}}$ is the identity protocol (no entangling unitary applied).
\end{conjecture}

\textbf{Why this conjecture is plausible.}
When $P = P_{\mathrm{id}}$, no entanglement is \emph{generated} by the protocol.
The witness $\mathcal{W}$ then measures the entanglement already present in the initial
state under the allowed operations.  For a Gibbs state at finite temperature,
Conjecture~\ref{conj:thermal} amounts to asking whether thermal equilibrium states
have operationally accessible entanglement (in the Wiseman-Vaccaro sense) under
realistic device operations.

For BCS superconductors specifically, there is a further suppression from the
charge superselection rule: under strict U(1) charge conservation, Cooper-pair
entanglement across a spatial cut is not operationally accessible
(Verstraete-Cirac, 2003~\cite{Verstraete2003}).  This provides partial
support for the conjecture in the superconducting case.

\textbf{Proof strategy for OP-new.}  Three layers of evidence are needed:
\begin{enumerate}[noitemsep]
  \item \textit{Free-fermion model.}  Compute $S_{\mathrm{cert}}(\rho_\beta)$
    for a 1D free-fermion chain at temperature $T$ under the constraint that
    the protocol $P$ consists only of single-site unitaries (no entangling gates).
    Prediction: $S_{\mathrm{cert}} = 0$ because no entanglement can be generated
    by single-site operations from a thermal product state.
  \item \textit{BCS-like model.}  Compute $S_{\mathrm{cert}}(\rho_{\mathrm{BCS}})$
    for the reduced state of a spatial half-chain under the charge SSR.
    Prediction: $S_{\mathrm{cert}} \leq E_{\mathcal{P}}^{\mathrm{SSR}}$ (accessible
    entanglement under SSR), which vanishes in the strict-SSR limit~\cite{Wiseman2003}.
  \item \textit{General argument.}  Use the monotonicity of relative entropy under
    CPTP maps to show that the accessible entanglement of $\rho_\beta$ under
    the device's operational restrictions is bounded by a quantity that vanishes
    as $T\to\infty$ and as the SSR becomes strict.
\end{enumerate}

A negative resolution of Conjecture~\ref{conj:thermal} --- showing that equilibrium
$S_{\mathrm{cert}}$ is large under realistic operations --- would require either
$E_* \ll E_P$ (rendering all near-term experiments insensitive) or a fundamental
reformulation of the sourcing mechanism.

\subsection{Status of $E_*$ and the constrained combination}
\label{sec:Estar}

The reference energy $E_*$ is an EFT matching coefficient that the framework does not
determine.  Only the product
\begin{equation}
  \Lambda_{\mathrm{eff}} \equiv \tilde{\kappa}\,E_*\,\alpha_{\mathrm{screen}}
  \quad [\mathrm{J/bit}]
\end{equation}
is experimentally constrained.  The three physically motivated choices span many orders
of magnitude:

\begin{center}\small
\begin{tabular}{@{}p{4.5cm}p{2.5cm}p{4cm}@{}}
\toprule
Choice & $E_*$ & $\Delta a$ (100 bits, 100\,$\mu$m, $\alpha_{\mathrm{screen}}=1$) \\
\midrule
Planck scale ($E_* = E_P$) & $1.96\times10^9$\,J & $3.6\times10^{-9}$\,m/s$^2$ \\
Qubit energy ($E_* = \hbar\omega$, 5\,GHz) & $3.3\times10^{-24}$\,J & $6\times10^{-42}$\,m/s$^2$ \\
Thermal ($E_* = k_BT$, 15\,mK) & $2.1\times10^{-25}$\,J & $4\times10^{-43}$\,m/s$^2$ \\
\bottomrule
\end{tabular}
\end{center}

Only the Planck-scale choice produces a signal above any conceivable near-term threshold,
making it the benchmark hypothesis.  This sharpens the importance of OP-new (does thermal
$S_{\mathrm{cert}}$ vanish?), OP3 (what is $\alpha_{\mathrm{screen}}$ for realistic NISQ
circuits?), and OP10 (does $\tilde{\kappa}E_*$ renormalize from the Planck scale to the
laboratory?).


\section{Conclusions}

The entanglement--gravity framework has reached a state of internal consistency.
Three points define its honest status:

\textbf{1. The source is operationally defined.}
The gravitational source field $S_{\mathrm{cert}}(x)$ is protocol-certified
entanglement density: entanglement generated by a specified control sequence,
certified by an entanglement witness relative to a pre-registered bipartition,
and varied between experimental branches.  Background many-body correlations
--- including BCS Cooper-pair entanglement and all thermal correlations ---
are excluded by the Cancellation Proposition~\ref{prop:cancel}: they are
branch-independent and cancel exactly in the differential observable.

\textbf{2. The observable is $\Lambda_{\mathrm{eff}} = \tilde{\kappa}E_*\alpha_{\mathrm{screen}}$.}
Neither $\tilde{\kappa}$, $E_*$, nor $\alpha_{\mathrm{screen}}$ is individually constrained
by any experiment.  Only their product is.  Mapping a null result to a bound on $\tilde{\kappa}$
requires independent input from OP3 (Lindblad numerics for $\alpha_{\mathrm{screen}}$) and
OP10 (RG flow of $\tilde{\kappa}E_*$ from UV to laboratory scales).

\textbf{3. Force sensing is the correct experimental path.}
For $E_* = E_P$, $|\tilde{\kappa}| = 1/4$, $\Delta S_{\mathrm{cert}} = 100$\,bits,
$R = 100\,\mu$m, the predicted unsuppressed acceleration is
$\Delta a = 3.6\times10^{-9}$\,m/s$^2$.  Cryogenic SQUID-based force sensors
achieve $\sim10^{-10}$\,m/s$^2/\sqrt{\mathrm{Hz}}$, placing the signal above the
noise floor by a factor of $\sim\!36$.  Ramsey-interferometry phase shifts are
suppressed by $c^2$ and lie far below current thresholds.

The framework is now \emph{falsifiable and tightly constrained}: a well-designed
differential force-sensing experiment, combined with theoretical progress on
OP-new (Conjecture~\ref{conj:thermal}), OP3, and OP10, will either detect an
entanglement-dependent gravitational anomaly or rule it out at the Planck-scale
benchmark level.

\section*{Acknowledgments}
The author thanks AI research assistants for literature synthesis, dimensional
verification, and critical review.  All physical assumptions, modeling choices,
and scientific judgments are the author's own.

\begin{thebibliography}{99}

\bibitem{Jacobson1995}
T.~Jacobson, \textit{Phys.\ Rev.\ Lett.}\ \textbf{75}, 1260 (1995).

\bibitem{Verlinde2017}
E.~Verlinde, \textit{SciPost Phys.}\ \textbf{2}, 016 (2017).

\bibitem{Wiseman2003}
H.~M.~Wiseman and J.~A.~Vaccaro,
``Entanglement of indistinguishable particles shared between two parties,''
\textit{Phys.\ Rev.\ Lett.}\ \textbf{91}, 097903 (2003).

\bibitem{Verstraete2003}
F.~Verstraete and J.~I.~Cirac,
``Quantum nonlocality in the presence of superselection rules and data hiding protocols,''
\textit{Phys.\ Rev.\ Lett.}\ \textbf{91}, 010404 (2003).

\bibitem{Bartlett2007}
S.~D.~Bartlett, T.~Rudolph, and R.~W.~Spekkens,
``Reference frames, superselection rules, and quantum information,''
\textit{Rev.\ Mod.\ Phys.}\ \textbf{79}, 555 (2007).

\bibitem{Asenbaum2020}
P.~Asenbaum \textit{et al.}, \textit{Phys.\ Rev.\ Lett.}\ \textbf{125}, 191101 (2020).

\bibitem{MICROSCOPE2022}
P.~T.~Touboul \textit{et al.}, \textit{Phys.\ Rev.\ Lett.}\ \textbf{129}, 121102 (2022).

\bibitem{Blume2017}
R.~Blume-Kohout \textit{et al.}, \textit{Nat.\ Commun.}\ \textbf{8}, 14485 (2017).

\end{thebibliography}

\end{document}
